% =====================================================================
% CHAPTER 1: Introduction (new department template, 2026-07-14)
% Red sections fixed: Overview, Background and Motivation, Research
% Questions / Problem Statements, Research Objectives, Key Contributions,
% Organization of the Report. No Summary section.
% Assembled per thesis/RESTRUCTURE_PLAN_V2.md (ch1 owner) from:
%   - old mainmatter/ch1.tex 1.1 Overview, 1.2 Background/Motivation,
%     1.3 Objectives, 1.7 Summary (dissolved).
%   - old mainmatter/ch3.tex 3.2 Problem Statements (recast as RQs).
%   - old mainmatter/ch9.tex contributions + frontmatter/abstract.tex
%     (contribution grounding; numbers frozen with sources).
% All numbers copied verbatim from the sources above; none introduced.
% =====================================================================
\chapter{Introduction}
\label{ch:intro}

\section{Overview}
A machine that must act in three-dimensional space first needs to know
how far away things are. A mobile robot needs it to avoid obstacles, a
drone needs it to hold position, and an augmented-reality headset needs
it to anchor virtual content to real surfaces. Binocular stereo recovers
this depth passively from two inexpensive image sensors, but the dense
correspondence it requires is computationally demanding, and the deep
networks that now solve it most accurately are far too large and too slow
for the embedded accelerators these platforms carry. This thesis designs,
trains, and deploys a compact stereo matching network that narrows this
gap. The network produces dense disparity, and from it metric
three-dimensional reconstruction, within the parameter, latency, and
memory budget of edge hardware.

The problem it addresses is a joint one: dense and accurate stereo depth
on hardware that supplies only a few tera-operations per second, a few
gigabytes of memory, and a few watts of power. This chapter states the
background and motivation, the research questions the thesis answers, the
objectives, the scope of the work, the key contributions, and the
organization of the report.

\FloatBarrier
\section{Background and Motivation}
Depth is not an optional cue. It is the quantity on which the safety and
usefulness of spatially acting systems depend, and the sensing options
that supply it differ sharply in cost and character. Scanning LiDAR
measures time of flight along emitted laser beams and returns accurate but
sparse points, at a price and power draw that exclude small platforms.
Structured-light and time-of-flight cameras emit infrared patterns or
pulses and recover dense depth indoors, but their emitters limit range,
fail in sunlight, and interfere with one another when several devices
share a space. Among the passive ways to recover depth, binocular stereo
suits embedded systems because it emits nothing, draws little power, works
indoors and outdoors alike, and uses two inexpensive image sensors rather
than a laser scanner. Its cost is computational. The correspondence
between the two views must be solved for every pixel, and that computation
is what this thesis sets out to make affordable.

Binocular stereo recovers depth from the horizontal shift of a scene
point between two horizontally separated views. That shift, the
disparity, is inversely proportional to depth through a calibrated
pinhole relation. A dense disparity map
therefore yields a dense depth map, which downstream perception and
reconstruction consume directly, and computing that disparity map
accurately and quickly is the task this thesis addresses.

The systems considered in this thesis operate under similar constraints:
limited computational resources, restricted power availability, and a
need for real-time estimates. Their depth requirements, however, are not
identical. A mobile robot primarily requires dependable obstacle depth
over short and medium ranges, whereas a drone also demands a high update
rate because it may approach obstacles more rapidly. In augmented
reality, accurate object boundaries are especially important for
realistic occlusion between virtual and physical content. A compact
sensing platform, by contrast, places greater emphasis on metric accuracy
for measurement and three-dimensional reconstruction. An embedded stereo
network intended for these settings must therefore achieve a practical
compromise among accuracy, boundary fidelity, inference time, and memory
use. Establishing this compromise is a central concern of the proposed
design.

\nm{$d$}{disparity, the horizontal pixel shift between rectified views}
\nm{$Z$}{depth of a scene point along the camera axis}
\nm{$f$}{camera focal length in pixels}
\nm{$B$}{stereo baseline, the distance between the two camera centers}
\nm{$(c_x, c_y)$}{principal point of the camera}
\nm{EPE}{end-point error, the mean absolute disparity error}
\nm{bad-$t$}{percentage of valid pixels with error above $t$ pixels}
\nm{D1-all}{percentage of pixels with error above 3 px and 5\% of ground truth}
\nm{GRU}{gated recurrent unit}
\nm{GEV}{geometry encoding volume}
\nm{KD}{knowledge distillation}
\nm{SGM}{semi-global matching}
\nm{FPS}{frames per second}
\nm{TOPS}{tera-operations per second}
\nm{INT8}{eight-bit integer arithmetic}
\nm{GT}{ground truth}

Deep networks have transformed the accuracy of stereo matching. The most
accurate systems now pair a foundation-scale monocular prior with an
iterative refinement core; FoundationStereo, for example, reports 0.34
pixel end-point error on the SceneFlow benchmark with a backbone of
roughly 340 million parameters \cite{wen2025foundationstereo}. That
accuracy comes at a computational cost that embedded hardware cannot pay.
The edge accelerators found on mobile robots and drones deliver a few
tera-operations per second, carry a few gigabytes of memory, and operate
within a power envelope of a few watts to a few tens of watts. A network
with hundreds of millions of parameters and hundreds of milliseconds of
latency does not fit that budget. The problem this thesis addresses is
therefore not accuracy alone but the joint objective of accuracy, latency,
and memory: a stereo network that is accurate enough to be useful, small
enough to fit in on-device memory, and fast enough to run in real time on
an edge accelerator. Figure \ref{fig:edge_gap} places the published
literature on this trade-off and marks the region this thesis targets.

\begin{figure}[!htbp]
\centering
\includegraphics[width=\textwidth]{fig_1_3_edge_gap}
\caption{Trainable parameter count against SceneFlow accuracy, with the targeted edge budget shaded.}
\label{fig:edge_gap}
\end{figure}

\FloatBarrier
\section{Research Questions / Problem Statements}
Stereo matching takes a pair of rectified images and produces a dense
disparity map, from which metric depth follows directly through the
calibrated relation between disparity and depth. This thesis addresses
that task under a hardware budget rather than an accuracy budget alone:
a few million parameters, real-time latency on the target device, a
small inference working set, and integer-friendly operators. The design
assumes rectified binocular input, a maximum disparity of 192 pixels at
the inference resolution, and training on synthetic data with dense
ground truth.

Those requirements frame four research questions, each of which the thesis
answers with measured evidence:
\begin{enumerate}
\item \textbf{RQ1.} Can a stereo network within a few million parameters,
and inside the sub-3 million edge budget, reach competitive in-domain
accuracy when trained on synthetic data with dense ground truth?
\item \textbf{RQ2.} Does such a compact, synthetic-only network generalize
zero-shot to an unseen real domain, without any target-domain adaptation?
\item \textbf{RQ3.} Can it meet the real-time latency target on the intended
edge device class, after inference optimization and eight-bit deployment?
\item \textbf{RQ4.} Does the design tolerate the small vertical
misalignments of imperfect rectification, the condition a low-cost
stereo rig produces, which is the second objective of this thesis?
\end{enumerate}

\FloatBarrier
\section{Research Objectives}
The thesis is built around two objectives, one for each half of the
problem set out above: making stereo matching cheap enough to run on
edge hardware, and making it robust to the imperfect camera geometry
that such hardware supplies in practice. The objectives of this thesis
are:
\begin{enumerate}
\item To design a computationally efficient stereo matching pipeline,
leveraging AI-based disparity refinement to enhance depth estimation on
resource-limited platforms.
\item To design an architecture able to withstand camera rectification
imperfections.
\end{enumerate}

The first objective is the heart of the thesis. Stereo matching is
accurate today mainly because the leading networks are very large; the
aim here is to keep most of that accuracy while shedding the size, so
that depth estimation can run on the small, low-power computers that
robots, drones, headsets, and handheld sensors actually carry, machines
that offer only a few tera-operations per second, a few gigabytes of
memory, and a few watts of power. The phrase \emph{AI-based disparity
refinement} in the objective points to how this is done. The largest
networks build one big cost volume and read the answer off it in a
single step, and that step is where most of their cost lies. StereoLite
instead starts from a rough disparity for each small image region and
improves it a few times with a learned update; it is this repeated
refinement, rather than a single expensive pass, that recovers accuracy
without a large parameter count. Chapter~\ref{ch:design} weighs the
choice against the alternatives in full. The objective is met when the
first three research questions are answered together: accuracy in
domain, accuracy on an unseen real domain, and real-time speed on the
edge board once the model is optimised and quantised.

The second objective addresses something the first does not. A network
can score well on benchmark image pairs, which are aligned almost
perfectly, and still struggle on a real camera, because a cheap stereo
rig is never perfectly aligned: its two views are typically off by a
fraction of a pixel to about a pixel in the vertical direction. An
architecture that only works on flawless input is not yet ready for the
hardware it is built for, so the second objective asks the network to
hold its accuracy when that small misalignment is present. The thesis
checks this two ways, which together answer the fourth research
question: by shifting one image vertically by a known amount and
watching how much the error grows, and by running the trained model on
the project's own, imperfectly aligned rig to confirm the error stays
low.

One might ask why this thesis designs a new compact network instead of
compressing an existing large one. There are two reasons. The accuracy of
foundation-scale systems rests largely on their monocular priors and
their sheer scale, and neither survives aggressive compression, so
compressing such systems tends to remove the very component that makes
them accurate. In addition, the tile-hypothesis and recurrent-refinement
families were built for efficiency from the start and give a cleaner
starting point for an edge budget.

\FloatBarrier
\section{Scope of the Research}
\label{sec:scope}
It helps to state plainly what this thesis does and does not attempt, so
the results are read against the right yardstick. The work is limited to
passive, rectified, binocular stereo: depth is recovered from two
horizontally separated cameras whose images have already been rectified,
and nothing is emitted into the scene. The network is trained under
supervision on synthetic image pairs with dense ground truth, kept
compact at a few million parameters, and evaluated in three ways, on the
synthetic test set it shares a domain with, zero-shot on an unseen real
benchmark, and on the project's own stereo rig. Its inference is then
optimized and quantized to eight bits so it can be timed on the intended
edge board rather than only on a desktop GPU. The recovered disparity is
finally turned into metric depth and, for the rig scenes, into colored
three-dimensional point clouds through the calibrated camera geometry;
this reconstruction is the one named in the thesis title.

Several neighboring problems are deliberately left out, not because they
lack value but because each would pull the work away from the edge-stereo
question at its center. The thesis does not use active depth sensors such
as LiDAR or structured light, does not train without labels through
self-supervision, and does not estimate depth from a single image. It
also stops at a working, measured research prototype: packaging the model
as a deployment software development kit, and the systems engineering that
a shipped product would require, are outside its scope. Where these
boundaries touch a result, the relevant chapter says so, and the ones
worth revisiting are collected as future work in
Chapter~\ref{ch:conclusion}.

\FloatBarrier
\section{Key Contributions}
Meeting the two objectives above required building something (a
network), deciding how to build it defensibly (a methodology), and then
showing with measurements that both hold up (a set of results). The
contributions of this thesis fall into those three kinds.

\textbf{An architecture.} This thesis contributes StereoLite, a 2.96
million parameter edge-budget stereo network purpose-built for objective
1. It combines a truncated YOLO feature encoder, a left-image recurrent
context stream, plane-aware coarse-to-fine tile refinement, and a
fail-soft quarter-resolution geometry encoding volume, four blocks
chosen and combined specifically to keep parameter count, latency, and
memory inside the edge budget without discarding the accuracy that
larger cost-volume networks get from full-resolution search.

\textbf{A design methodology.} Building that network raised a recurring
question: how does a designer justify keeping one block over an
alternative under a hard parameter budget, rather than by intuition
alone? This thesis answers it with an evidence-gated design methodology
that admits an architectural choice only after it wins a controlled
single-knob ablation on a fixed 100-pair harness, where competing
variants share data, seed, and step budget so a measured difference is
attributable to the design change alone. Every block listed above
passed this test before being kept. The same evidence-gathering habit
also produced a survey of the efficient and edge-deployable stereo
families that grounds the design choices, stating what edge-deployable
stereo requires and mapping the accuracy, latency, and memory
trade-offs of the published literature.

\textbf{Results that follow from applying both.} Four measured outcomes
demonstrate that the architecture and the methodology deliver on the
two objectives:
\begin{itemize}
\item \emph{In-domain accuracy (objective 1).} The trained checkpoint
reaches 0.78 pixel end-point error and 3.40 percent D1-all on the full
FlyingThings3D test set of 4,370 pairs, while running in real time on a
laptop RTX 3050 GPU at 49.8 ms per forward pass at $384 \times 640$ in
half precision.
\item \emph{Cross-domain accuracy (objective 1).} The same training
pipeline gives the compact network 10.9 percent zero-shot Middlebury
2014 D1-all on a domain it never trains on, within four points of the
larger distillation-trained LiteAnyStereo.
\item \emph{Edge deployment (objective 1).} An equivalence-proven
inference-optimization pass cuts latency by a factor of 1.74 with no
loss of accuracy and identifies the operators that block eight-bit
deployment, so that the optimized eight-bit model runs in 36.3 ms, or
27.5 frames per second, on the Jetson Orin Nano.
\item \emph{Rectification tolerance (objective 2).} Under a controlled
vertical misalignment sweep, end-point error rises only from 1.03 to
1.53 px across the sub-pixel to one-pixel range a calibrated rig
produces, and on the project's own rig the predicted disparity
back-projects into metric colored point clouds on real indoor scenes.
\end{itemize}

\FloatBarrier
\section{Organization of the Report}
The rest of this report is arranged as follows.
\begin{itemize}
\item Chapter~\ref{ch:lit} surveys the stereo-matching literature, moving
from the classical pipeline through the deep-learning eras to the efficient
and foundation-model families. It states what edge-deployable stereo
requires and compares the published methods on accuracy, latency, and memory.
\item Chapter~\ref{ch:design} sets out the engineering design and
methodology of StereoLite. It covers the design requirements and criteria,
the design steps, the controlled ablation methodology, and the
block-by-block architecture together with its assessment metrics.
\item Chapter~\ref{ch:implementation} describes the system development and
implementation: the selection of parameters and components, the training
objective and protocol, the software realization, and the edge-inference
optimization.
\item Chapter~\ref{ch:experimental} covers the experimental setup and data
collection, including the datasets, the stereo rig and calibration, the test
benches, and the measurement procedures.
\item Chapter~\ref{ch:results} presents and discusses the results: the
training dynamics, the ablation studies, the in-domain and zero-shot
accuracy, the rectification robustness, the efficiency and deployment
measurements, and the three-dimensional reconstruction.
\item Chapter~\ref{ch:socio} considers the societal, health, and cultural
impact of the work, its safety and ethical dimensions, and its alignment
with the Sustainable Development Goals, including its environmental and
sustainability case.
\item Chapter~\ref{ch:management} records how the thesis was managed, its
time frame, its overall and component-level budget, its resource planning,
and the division of work between the authors.
\item Chapter~\ref{ch:conclusion} concludes the report, states its
limitations, and outlines future work.
\end{itemize}
The appendices open with the modern tools used in the work (Appendix A)
and the program-outcome mapping (Appendix B), then address the knowledge
profiles, the complex engineering problems, and the complex engineering
activities engaged by the thesis (Appendices C to E), and close with the
technical specifications, the plagiarism report, and the AI writing
report (Appendices F to H).

% SOURCES: old mainmatter/ch1.tex 1.1 Overview, 1.2 Background and
%   Motivation (with nm{} nomenclature block + fig:edge_gap intact), 1.3
%   Objectives (-> 1.4 Research Objectives), 1.7 Summary (dissolved into
%   1.1 and 1.6). old mainmatter/ch3.tex 3.2 Problem Statements + the
%   design-assumptions line and envelope summary from 3.3 (recast as 1.3
%   RQ1-RQ4 + problem-statement prose). old mainmatter/ch9.tex contribution
%   itemize + frontmatter/abstract.tex (1.5 Key Contributions grounding;
%   all numbers copied verbatim: 2.96 M params, 0.78 px EPE, 3.40% D1-all,
%   4,370 pairs, 10.9% MB14 zero-shot D1-all, 49.8 ms fp16 RTX 3050,
%   36.3 ms / 27.5 FPS Orin Nano, 1.74x latency cut, 1.03-1.53 px
%   rectification sweep, within four points of LiteAnyStereo).
% DROPPED: old 1.1 sentence listing "research methodology" and "plan" as
%   chapter contents (those sections moved away; replaced with the new
%   section list). Verbose 1.7 Summary recap (dissolved, no Summary section
%   per template).
% MOVED-AWAY: old ch1 1.4 Research Methodology + its 6-step enumerate ->
%   ch3 (ch:design), owner ch3. old ch1 1.5 Project Planning + Work Plan +
%   fig:gantt (fig_b_1_gantt_fixed) -> ch8 8.2.1 (ch:management), owner ch8.
%   Chapter-2 depth/disparity equation eq:depth_from_disparity lives in
%   ch:lit (owner ch2); referenced here in prose only, not via \eqref.
